\documentclass[10pt,landscape]{article}
% UTF-8 Schriftkodierung
\usepackage[utf8]{inputenc}\usepackage[T1]{fontenc}
% get rid of the "no room for new \ldots". Standard TeX registers are limitied to 256, etex increases the max.
\usepackage{etex}
\usepackage{multicol}
\usepackage{calc}
\usepackage{ifthen}
\usepackage[landscape]{geometry}
\usepackage{aass-common}
\usepackage{aass-units}

%%%%%%%%%%%%%%%%%%%%%%%%%%%%%%%%%%%%%%%%%%%%%%%%%%%%%%%%%%%
%%% Colors
\usepackage[svgnames,hyperref]{xcolor}
\input{./defs/colors.def}

% To make this come out properly in landscape mode, do one of the following
% 1.
%  pdflatex latexsheet.tex
%
% 2.
%  latex latexsheet.tex
%  dvips -P pdf  -t landscape latexsheet.dvi
%  ps2pdf latexsheet.ps


% If you're reading this, be prepared for confusion.  Making this was
% a learning experience for me, and it shows.  Much of the placement
% was hacked in; if you make it better, let me know...


% 2008-04
% Changed page margin code to use the geometry package. Also added code for
% conditional page margins, depending on paper size. Thanks to Uwe Ziegenhagen
% for the suggestions.

% 2006-08
% Made changes based on suggestions from Gene Cooperman. <gene at ccs.neu.edu>


% To Do:
% \listoffigures \listoftables
% \setcounter{secnumdepth}{0}


% This sets page margins to .5 inch if using letter paper, and to 1cm
% if using A4 paper. (This probably isn't strictly necessary.)
% If using another size paper, use default 1cm margins.
\ifthenelse{\lengthtest { \paperwidth = 11in}}
	{ \geometry{top=.5in,left=.5in,right=.5in,bottom=.5in} }
	{\ifthenelse{ \lengthtest{ \paperwidth = 297mm}}
		{\geometry{top=1cm,left=1cm,right=1cm,bottom=1cm} }
		{\geometry{top=1cm,left=1cm,right=1cm,bottom=1cm} }
	}

% Turn off header and footer
\pagestyle{empty}
 

% Redefine section commands to use less space
\makeatletter
\renewcommand{\section}{\@startsection{section}{1}{0mm}%
                                {-1ex plus -.5ex minus -.2ex}%
                                {0.5ex plus .2ex}%x
                                {\normalfont\large\bfseries}}
\renewcommand{\subsection}{\@startsection{subsection}{2}{0mm}%
                                {-1explus -.5ex minus -.2ex}%
                                {0.5ex plus .2ex}%
                                {\normalfont\normalsize\bfseries}}
\renewcommand{\subsubsection}{\@startsection{subsubsection}{3}{0mm}%
                                {-1ex plus -.5ex minus -.2ex}%
                                {1ex plus .2ex}%
                                {\normalfont\small\bfseries}}
\makeatother

% Define BibTeX command
\def\BibTeX{{\rm B\kern-.05em{\sc i\kern-.025em b}\kern-.08em
    T\kern-.1667em\lower.7ex\hbox{E}\kern-.125emX}}

% Don't print section numbers
\setcounter{secnumdepth}{0}


\setlength{\parindent}{0pt}
\setlength{\parskip}{0pt plus 0.5ex}


% -----------------------------------------------------------------------




\usepackage{listings}

\begin{document}




\lstset{% general command to set parameter(s)
language=[LaTeX]TeX,
%language=C,
extendedchars=true,
inputencoding=utf8,
breaklines=true;
basicstyle=\ttfamily\small, % print whole listing small
keywordstyle=\color{ubuntured}\bfseries\underbar,
% underlined bold black keywords
identifierstyle=\color{darkBlue}\bfseries, % nothing happens
commentstyle=\color{Grey}, % white comments
stringstyle=\ttfamily, % typewriter type for strings
texcsstyle=\color{ubuntured},
moretexcs={textcolor,itemhook},
emph={
A,
AuthorNotes,
Ca,
Cave,
Definition,
E,
Es,
Fazit,
H,
Layout,
Married,
MarriedNG,
MarriedLI,
Massnahmen,
Parallel,
T,
Z},
emphstyle=\color{Terminus}\bfseries,
%directivestyle=color{Red},
showstringspaces=false,
%morecomment=[s][\color{SandyBrown}]{\{}{\}},
%moredelim=[s][\color{DarkViolet}]{\{}{\}},
%morecomment=[l][\color{Grey}]{\%}
}





\raggedright
\footnotesize
\begin{multicols}{3}


% multicol parameters
% These lengths are set only within the two main columns
%\setlength{\columnseprule}{0.25pt}
\setlength{\premulticols}{1pt}
\setlength{\postmulticols}{1pt}
\setlength{\multicolsep}{1pt}
\setlength{\columnsep}{2pt}

\begin{center}
     \Large{\textbf{\LaTeXe\ Cheat Sheet}} \\
\end{center}

% \section{Document classes}
% \begin{tabular}{@{}ll@{}}
% \verb!book!    & Default is two-sided. \\
% \verb!report!  & No \verb!\part! divisions. \\
% \verb!article! & No \verb!\part! or \verb!\chapter! divisions. \\
% \verb!letter!  & Letter (?). \\
% \verb!slides!  & Large sans-serif font.
% \end{tabular}
% 
% Used at the very beginning of a document:
% \verb!\documentclass{!\textit{class}\verb!}!.  Use
% \verb!\begin{document}! to start contents and \verb!\end{document}! to
% end the document.
% 
% 
% \subsection{Common \texttt{documentclass} options}
\newlength{\MyLen}
\settowidth{\MyLen}{\texttt{letterpaper}/\texttt{a4paper} \ }
% \begin{tabular}{@{}p{\the\MyLen}%
%                 @{}p{\linewidth-\the\MyLen}@{}}
% \texttt{10pt}/\texttt{11pt}/\texttt{12pt} & Font size. \\
% \texttt{letterpaper}/\texttt{a4paper} & Paper size. \\
% \texttt{twocolumn} & Use two columns. \\
% \texttt{twoside}   & Set margins for two-sided. \\
% \texttt{landscape} & Landscape orientation.  Must use
%                      \texttt{dvips -t landscape}. \\
% \texttt{draft}     & Double-space lines.
% \end{tabular}
% 
% Usage:
% \verb!\documentclass[!\textit{opt,opt}\verb!]{!\textit{class}\verb!}!.
% 
% 
% \subsection{Packages}
% \settowidth{\MyLen}{\texttt{multicol} }
% \begin{tabular}{@{}p{\the\MyLen}%
%                 @{}p{\linewidth-\the\MyLen}@{}}
% %\begin{tabular}{@{}ll@{}}
% \texttt{fullpage}  &  Use 1 inch margins. \\
% \texttt{anysize}   &  Set margins: \verb!\marginsize{!\textit{l}%
%                         \verb!}{!\textit{r}\verb!}{!\textit{t}%
%                         \verb!}{!\textit{b}\verb!}!.            \\
% \texttt{multicol}  &  Use $n$ columns: 
%                         \verb!\begin{multicols}{!$n$\verb!}!.   \\
% \texttt{latexsym}  &  Use \LaTeX\ symbol font. \\
% \texttt{graphicx}  &  Show image:
%                         \verb!\includegraphics[width=!%
%                         \textit{x}\verb!]{!%
%                         \textit{file}\verb!}!. \\
% \texttt{url}       & Insert URL: \verb!\url{!%
%                         \textit{http://\ldots}%
%                         \verb!}!.
% \end{tabular}
% 
% Use before \verb!\begin{document}!. 
% Usage: \verb!\usepackage{!\textit{package}\verb!}!
% 
% 
% \subsection{Title}
% \settowidth{\MyLen}{\texttt{.author.text.} }
% \begin{tabular}{@{}p{\the\MyLen}%
%                 @{}p{\linewidth-\the\MyLen}@{}}
% \verb!\author{!\textit{text}\verb!}! & Author of document. \\
% \verb!\title{!\textit{text}\verb!}!  & Title of document. \\
% \verb!\date{!\textit{text}\verb!}!   & Date. \\
% \end{tabular}
% 
% These commands go before \verb!\begin{document}!.  The declaration
% \verb!\maketitle! goes at the top of the document.
% 
% \subsection{Miscellaneous}
% \settowidth{\MyLen}{\texttt{.pagestyle.empty.} }
% \begin{tabular}{@{}p{\the\MyLen}%
%                 @{}p{\linewidth-\the\MyLen}@{}}
% \verb!\pagestyle{empty}!     &   Empty header, footer
%                                  and no page numbers. \\
% 
% \end{tabular}

\section{Allgemein}

Leerzeilen markieren Absatzgrenzen.

{\LaTeX} kennt Befehle und Umgebungen. Befehlen können Argumente folgen, sie werden mit
geschwungenen Klammern \verb!{}! umschlossen, optionale
Befehle mit eckigen Klammern \verb![]!.

Beispiel: Befehle mit einem Argument, einem optionalen und normalen Argument und drei Argumenten

\verb!\Befehl{!\textit{Argument}\verb!}!\\
\verb!\Befehl[!\textit{Optionales Argument}\verb!]{!\textit{Argument}\verb!}!\\
\verb!\Befehl{!\textit{Argument 1}\verb!}{!\textit{Argument 2}\verb!}{!\textit{Argument 3}\verb!}!

Umgebungen werden mittels \verb!\begin{!\textit{Umgebungsname}\verb!}! begonnen und mittels \verb!\end{!\textit{Umgebungsname}\verb!}!  

\begin{verbatim}
\begin{center}
Dieser Text steht in der Umgebung "center" 
und wird daher zentriert.
\end{center}
\end{verbatim}

\begin{center}
Dieser Text steht in der Umgebung\\
"center" und wird daher zentriert.
\end{center}

Umgebungen können auch Argumente besitzen, z.\,B. \verb!\begin{multicols}{2}!.

\section{Gliederung}
\begin{multicols}{2}
\verb!\part{!\textit{title}\verb!}!  \\
\verb!\chapter{!\textit{title}\verb!}!  \\
\verb!\section{!\textit{title}\verb!}!  \\
\verb!\subsection{!\textit{title}\verb!}!  \\
\verb!\subsubsection{!\textit{title}\verb!}!  \\
\verb!\paragraph{!\textit{title}\verb!}!  \\
\verb!\subparagraph{!\textit{title}\verb!}!
\end{multicols}
{\raggedright
Die Gliederungskommandos können von einem \texttt{*}
gefolgt werden (z.\,B.
\verb!\section*{!\textit{title}\verb!}!) um eine
Nummerierung (und Eintrag ins Inhaltsverzeichnis) zu
verhindern. }

\section{Text}

\subsection{Textformatierungen}
\newcommand{\FontCmd}[3]{\PBS\verb!\#1{!\textit{text}\verb!}!  \> %
                         \verb!{\#2 !\textit{text}\verb!}! \> %
                         \#1{#3}}
% \begin{tabular}{@{}l@{}l@{}l@{}}
% \textit{Command} & \textit{Declaration} & \textit{Effect} \\
% \verb!\textrm{!\textit{text}\verb!}!                    & %
%         \verb!{\rmfamily !\textit{text}\verb!}!               & %
%         \textrm{Roman family} \\
% \verb!\textsf{!\textit{text}\verb!}!                    & %
%         \verb!{\sffamily !\textit{text}\verb!}!               & %
%         \textsf{Sans serif family} \\
% \verb!\texttt{!\textit{text}\verb!}!                    & %
%         \verb!{\ttfamily !\textit{text}\verb!}!               & %
%         \texttt{Typewriter family} \\
% \verb!\textmd{!\textit{text}\verb!}!                    & %
%         \verb!{\mdseries !\textit{text}\verb!}!               & %
%         \textmd{Medium series} \\
% \verb!\textbf{!\textit{text}\verb!}!                    & %
%         \verb!{\bfseries !\textit{text}\verb!}!               & %
%         \textbf{Bold series} \\
% \verb!\textup{!\textit{text}\verb!}!                    & %
%         \verb!{\upshape !\textit{text}\verb!}!               & %
%         \textup{Upright shape} \\
% \verb!\textit{!\textit{text}\verb!}!                    & %
%         \verb!{\itshape !\textit{text}\verb!}!               & %
%         \textit{Italic shape} \\
% \verb!\textsl{!\textit{text}\verb!}!                    & %
%         \verb!{\slshape !\textit{text}\verb!}!               & %
%         \textsl{Slanted shape} \\
% \verb!\textsc{!\textit{text}\verb!}!                    & %
%         \verb!{\scshape !\textit{text}\verb!}!               & %
%         \textsc{Small Caps shape} \\
% \verb!\emph{!\textit{text}\verb!}!                      & %
%         \verb!{\em !\textit{text}\verb!}!               & %
%         \emph{Emphasized} \\
% \verb!\textnormal{!\textit{text}\verb!}!                & %
%         \verb!{\normalfont !\textit{text}\verb!}!       & %
%         \textnormal{Document font} \\
% \verb!\underline{!\textit{text}\verb!}!                 & %
%                                                         & %
%         \underline{Underline}
% \end{tabular}

\begin{tabular}{@{}l@{}l@{}l@{}}
\textit{Befehl}  & \textit{Effekt} & \textit{Bedeutung} \\
\verb!\E{!\textit{Text}\verb!}!  & \E{Text} & Schwache Betonung \\
\verb!\EE{!\textit{Text}\verb!}!  & \EE{Text} & Starke Betonung \\
\verb!\T{!\textit{Text}\verb!}!  & \T{Text} & Terminus \\
\verb!\I{!\textit{Text}\verb!}!  & \I{Text} & Indexeintrag, schwache Betonung \\
\verb!\II{!\textit{Text}\verb!}!  & \II{Text} & Indexeintrag, starke Betonung \\
\verb!\Ca{!\textit{Text}\verb!}!  & \Ca{Text} & Cave! \\
\verb!\Z{!\textit{Text}\verb!}!  & \Z{Text} & Zusatzinfo \\
\end{tabular}

% The command (t\textit{tt}t) form handles spacing better than the
% declaration (t{\itshape tt}t) form.

% \subsection{Font size}
% \setlength{\columnsep}{14pt} % Need to move columns apart a little
% \begin{multicols}{2}
% \begin{tabbing}
% \verb!\footnotesize!          \= \kill
% \verb!\tiny!                  \>  \tiny{tiny} \\
% \verb!\scriptsize!            \>  \scriptsize{scriptsize} \\
% \verb!\footnotesize!          \>  \footnotesize{footnotesize} \\
% \verb!\small!                 \>  \small{small} \\
% \verb!\normalsize!            \>  \normalsize{normalsize} \\
% \verb!\large!                 \>  \large{large} \\
% \verb!\Large!                 \=  \Large{Large} \\  % Tab hack for new column
% \verb!\LARGE!                 \>  \LARGE{LARGE} \\
% \verb!\huge!                  \>  \huge{huge} \\
% \verb!\Huge!                  \>  \Huge{Huge}
% \end{tabbing}
% \end{multicols}
% \setlength{\columnsep}{1pt} % Set column separation back
% 
% These are declarations and should be used in the form
% \verb!{\small! \ldots\verb!}!, or without braces to affect the entire
% document.


\subsection{Striche}
\begin{tabular}{@{}llll@{}}
\textit{Name} & \textit{Source} & \textit{Example} & \textit{Usage} \\
hyphen  & \verb!-!   & X-ray          & Innerh. v. Wörtern \\
en-dash & \verb!--!  & 1--5           &  \\
% em-dash & \verb!---! & Yes---or no?    & Punctuation.
\end{tabular}


\subsection{Textumgebungen}
\settowidth{\MyLen}{\texttt{.begin.quotation.}}
\begin{tabular}{@{}p{\the\MyLen}%
                @{}p{\linewidth-\the\MyLen}@{}}
\verb!\begin{comment}!    &  Comment block (not printed). \\
\verb!\begin{quote}!      &  Indented quotation block. \\
\verb!\begin{quotation}!  &  Like \texttt{quote} with indented paragraphs. \\
\verb!\begin{verse}!      &  Quotation block for verse.
\end{tabular}

\subsection{Lists}
\settowidth{\MyLen}{\texttt{.begin.description.}}
\begin{tabular}{@{}p{\the\MyLen}%
                @{}p{\linewidth-\the\MyLen}@{}}
\verb!\begin{enumerate}!        &  Nummerierte Liste \\
\verb!\begin{itemize}!          &  Punktliste \\
\verb!\begin{description}!      &  Beschreibungsliste \\
\verb!\item! \textit{text}      &  Aufzählungspunkt \\
\verb!\item[!\textit{x}\verb!]! \textit{text}
                                &  \textit{x} wird als Aufzählungszeichen verwendet, Terminus für Beschreibungslisten\\
\end{tabular}

\subsection{Besondere Befehle}
\begin{tabular}{@{}l@{}l@{}l@{}}
\textit{Befehl}  & \textit{Effekt} & \textit{Bedeutung} \\
\verb!\VonBis{!\textit{a}\verb!}{!\textit{b}\verb!}!  & \VonBis{a}{b} &  \\

\end{tabular}



\subsection{Fließende Sachen}
\settowidth{\MyLen}{\texttt{.begin.equation..place.}}

Floating Bodies werden automatisch platziert. 

\begin{tabular}{@{}p{\the\MyLen}%
                @{}p{\linewidth-\the\MyLen}@{}}
\verb!\begin{table}[!\textit{Platz}\verb!]!     &  Fließende Tabelle \\
\verb!\begin{figure}[!\textit{Platz}\verb!]!    &  Grafik \\
\verb!\begin{equation}[!\textit{Platz}\verb!]!  &  Gleichung \\
\verb!\Captionof{!\textit{Typ}\verb!}{!\textit{text}\verb!}!           &  Beschriftung \\
\end{tabular}

Mit \textit{Platz} kann die Positionierung beeinflusst werden:  \texttt{t}=oben (top),
\texttt{h}=möglichst hier, \texttt{H}=unter allen Umständen hier, \texttt{b}=unten (bottom), \texttt{p}=eigene Seite, \texttt{!}=Anweisung befolgen, auch wenn es nicht gut ausschaut.  

%---------------------------------------------------------------------------



% \subsection{Verbatim text}
% 
% \settowidth{\MyLen}{\texttt{.begin.verbatim..} }
% \begin{tabular}{@{}p{\the\MyLen}%
%                 @{}p{\linewidth-\the\MyLen}@{}}
% \verb@\begin{verbatim}@ & Verbatim environment. \\
% \verb@\begin{verbatim*}@ & Spaces are shown as \verb*@ @. \\
% \verb@\verb!text!@ & Text between the delimiting characters (in this case %
%                       `\texttt{!}') is verbatim.
% \end{tabular}


% \subsection{Justification}
% \begin{tabular}{@{}ll@{}}
% \textit{Environment}  &  \textit{Declaration}  \\
% \verb!\begin{center}!      & \verb!\centering!  \\
% \verb!\begin{flushleft}!  & \verb!\raggedright! \\
% \verb!\begin{flushright}! & \verb!\raggedleft!  \\
% \end{tabular}

% \subsection{Miscellaneous}
% \verb!\linespread{!$x$\verb!}! changes the line spacing by the
% multiplier $x$.
% 


\section{Verweise}
\subsection{Querverweise}
\settowidth{\MyLen}{\texttt{.pageref.marker..}}
\begin{tabular}{@{}p{\the\MyLen}%
                @{}p{\linewidth-\the\MyLen}@{}}
\verb!\label{!\textit{xy}\verb!}!   &  Allgemein: Erstellt einen Referenzpunkt \textit{xy}
                          often of the form \verb!\label{sec:item}!. \\
\verb!\label{!\textit{chp:xy}\verb!}!   &  Erstellt einen Referenzpunkt für ein Kapitel\\
\verb!\label{!\textit{fig:xy}\verb!}!   &  Erstellt einen Referenzpunkt für eine Grafik (figure)\\
\verb!\label{!\textit{tab:xy}\verb!}!   &  Erstellt einen Referenzpunkt für eine Tabelle\\
\verb!\label{!\textit{topic:xy}\verb!}!   &  Erstellt einen Referenzpunkt für ein allgemeines Thema\\
\verb!\prettyref{!\textit{xy}\verb!}!   &  Gibt Referenznummer von \textit{xy} aus (Gliederungsnummer, Abbildungsnummer, \ldots) \\
\verb!\pageref{!\textit{xy}\verb!}! &  Gibt Seitennummer von \textit{xy} aus \\
\verb!\FullPrettyRef{!\textit{xy}\verb!}! &  Gibt Referenznummer und Seitennummer von \textit{xy} aus \\
\verb!\footnote{!\textit{text}\verb!}!  &  Fußnote \\
\verb!\ZFootnote{!\textit{text}\verb!}!  &  "`graue"' Fußnote \\
\end{tabular}

Bei Beschriftungen soll \verb!\label{!\textit{xy}\verb!}! innerhalb des Beschirftungstextes stehen.

% \begin{description}
%   \item[{\protect\verb!\label{!\textit{marker}\protect\verb!}!]
% Allgemein: Erstellt einen Referenzpunkt \textit{marker}
% \item[\protect\verb!\label{!\textit{chp:marker}\protect\verb!}!]
% Erstellt einen Referenzpunkt für ein Kapitel
% \item[\protect\verb!\label{!\textit{fig:marker}\protect\verb!}!] Erstellt einen Referenzpunkt für eine Grafik (figure)\\
% \item[\protect\verb!\label{!\textit{tab:marker}\protect\verb!}!]     Erstellt einen Referenzpunkt für eine Tabelle\\
% \item[\protect\verb!\label{!\textit{topic:marker}\protect\verb!}!]     Erstellt einen Referenzpunkt für ein allgemeines Thema\\
% \item[\protect\verb!\prettyref{!\textit{marker}\protect\verb!}!]     Gibt Referenznummer von \textit{marker} aus (Gliederungsnummer, Abbildungsnummer, \ldots) \\
% \item[\protect\verb!\pageref{!\textit{marker}\protect\verb!}!]  Gibt Seitennummer von \textit{marker} aus \\
% \item[\protect\verb!\FullPrettyRef{!\textit{marker}\protect\verb!}!]  Gibt Referenznummer und Seitennummer von \textit{marker} aus \\
% \item[\protect\verb!\footnote{!\textit{text}\protect\verb!}!]   Fußnote \\
% \item[\protect\verb!\ZFootnote{!\textit{text}\protect\verb!}!]   "`graue"' Fußnote 
% \end{description}


\subsection{Literaturzitate}

\begin{tabular}{@{}p{\the\MyLen}@{}p{\linewidth-\the\MyLen}@{}}
\verb!\cite{!\textit{key}\verb!}!       &
        Zitiert \textit{key} \\
% \verb!\citeA{!\textit{key}\verb!}!      &
%         Full author list. (Watson and Crick) \\
% \verb!\citeN{!\textit{key}\verb!}!      &
%         Full author list and year. Watson and Crick (1953) \\
% \verb!\shortcite{!\textit{key}\verb!}!  &
%         Abbreviated author list and year. ? \\
% \verb!\shortciteA{!\textit{key}\verb!}! &
%         Abbreviated author list. ? \\
% \verb!\shortciteN{!\textit{key}\verb!}! &
%         Abbreviated author list and year. ? \\
% \verb!\citeyear{!\textit{key}\verb!}!   &
%         Cite year only. (1953) \\
\end{tabular}

% 
% \subsection{Citation types}
% \settowidth{\MyLen}{\texttt{.shortciteN.key..}}
% \begin{tabular}{@{}p{\the\MyLen}@{}p{\linewidth-\the\MyLen}@{}}
% \verb!\cite{!\textit{key}\verb!}!       &
%         Full author list and year. (Watson and Crick 1953) \\
% \verb!\citeA{!\textit{key}\verb!}!      &
%         Full author list. (Watson and Crick) \\
% \verb!\citeN{!\textit{key}\verb!}!      &
%         Full author list and year. Watson and Crick (1953) \\
% \verb!\shortcite{!\textit{key}\verb!}!  &
%         Abbreviated author list and year. ? \\
% \verb!\shortciteA{!\textit{key}\verb!}! &
%         Abbreviated author list. ? \\
% \verb!\shortciteN{!\textit{key}\verb!}! &
%         Abbreviated author list and year. ? \\
% \verb!\citeyear{!\textit{key}\verb!}!   &
%         Cite year only. (1953) \\
% \end{tabular}
% 
% All the above have an \texttt{NP} variant without parentheses;
% Ex. \verb!\citeNP!.


% \subsection{\BibTeX\ entry types}
% \settowidth{\MyLen}{\texttt{.mastersthesis.}}
% \begin{tabular}{@{}p{\the\MyLen}@{}p{\linewidth-\the\MyLen}@{}}
% \verb!@article!         &  Journal or magazine article. \\
% \verb!@book!            &  Book with publisher. \\
% \verb!@booklet!         &  Book without publisher. \\
% \verb!@conference!      &  Article in conference proceedings. \\
% \verb!@inbook!          &  A part of a book and/or range of pages. \\
% \verb!@incollection!    &  A part of book with its own title. \\
% %\verb!@manual!          &  Technical documentation. \\
% %\verb!@mastersthesis!   &  Master's thesis. \\
% \verb!@misc!            &  If nothing else fits. \\
% \verb!@phdthesis!       &  PhD. thesis. \\
% \verb!@proceedings!     &  Proceedings of a conference. \\
% \verb!@techreport!      &  Tech report, usually numbered in series. \\
% \verb!@unpublished!     &  Unpublished. \\
% \end{tabular}
% 
% \subsection{\BibTeX\ fields}
% \settowidth{\MyLen}{\texttt{organization.}}
% \begin{tabular}{@{}p{\the\MyLen}@{}p{\linewidth-\the\MyLen}@{}}
% \verb!address!         &  Address of publisher.  Not necessary for major
%                                 publishers.  \\
% \verb!author!           &  Names of authors, of format .... \\
% \verb!booktitle!        &  Title of book when part of it is cited. \\
% \verb!chapter!          &  Chapter or section number. \\
% \verb!edition!          &  Edition of a book. \\
% \verb!editor!           &  Names of editors. \\
% \verb!institution!      &  Sponsoring institution of tech.\ report. \\
% \verb!journal!          &  Journal name. \\
% \verb!key!              &  Used for cross ref.\ when no author. \\
% \verb!month!            &  Month published. Use 3-letter abbreviation. \\
% \verb!note!             &  Any additional information. \\
% \verb!number!           &  Number of journal or magazine. \\
% \verb!organization!     &  Organization that sponsors a conference. \\
% \verb!pages!            &  Page range (\verb!2,6,9--12!). \\
% \verb!publisher!        &  Publisher's name. \\
% \verb!school!           &  Name of school (for thesis). \\
% \verb!series!           &  Name of series of books. \\
% \verb!title!            &  Title of work. \\
% \verb!type!             &  Type of tech.\ report, ex. ``Research Note''. \\
% \verb!volume!           &  Volume of a journal or book. \\
% \verb!year!             &  Year of publication. \\
% \end{tabular}
% Not all fields need to be filled.  See example below.

% % \subsection{Common \BibTeX\ style files}
% % \begin{tabular}{@{}l@{\hspace{1em}}l@{\hspace{3em}}l@{\hspace{1em}}l@{}}
% % \verb!abbrv!    &  Standard &
% % \verb!abstract! &  \texttt{alpha} with abstract \\
% % \verb!alpha!    &  Standard &
% % \verb!apa!      &  APA \\
% % \verb!plain!    &  Standard &
% % \verb!unsrt!    &  Unsorted \\
% % \end{tabular}
% % 
% % The \LaTeX\ document should have the following two lines just before
% % \verb!\end{document}!, where \verb!bibfile.bib! is the name of the
% % \BibTeX\ file.
% % \begin{verbatim}
% % \bibliographystyle{plain}
% % \bibliography{bibfile}
% % \end{verbatim}
% % 
% % \subsection{\BibTeX\ example}
% % The \BibTeX\ database goes in a file called
% % \textit{file}\texttt{.bib}, which is processed with \verb!bibtex file!. 
% % \begin{verbatim}
% % @String{N = {Na\-ture}}
% % @Article{WC:1953,
% %   author  = {James Watson and Francis Crick},
% %   title   = {A structure for Deoxyribose Nucleic Acid},
% %   journal = N,
% %   volume  = {171},
% %   pages   = {737},
% %   year    = 1953
% % }
% % \end{verbatim}

\subsection{Kommentare}

\verb!\Commentary[!\textit{Titel}\verb!]{!\textit{Kommentartext}\verb!}!

Erstellt einen Eintrag in das Kommentarverzeichnis.

\section{Symbole}

% \subsection{Symbols}
\begin{tabular}{@{}l@{\hspace{1em}}l@{\hspace{2em}}l@{\hspace{1em}}l@{\hspace{2em}}l@{\hspace{1em}}l@{\hspace{2em}}l@{\hspace{1em}}l@{}}
\&              &  \verb!\&! &
\_              &  \verb!\_! &
\ldots          &  \verb!\ldots! &
\textbullet     &  \verb!\textbullet! \\
\$              &  \verb!\$! &
\^{}            &  \verb!\^{}! &
\textbar        &  \verb!\textbar! &
\textbackslash  &  \verb!\textbackslash! \\
\%              &  \verb!\%! &
\~{}            &  \verb!\~{}! &
\#              &  \verb!\#! &
\S              &  \verb!\S! \\
\end{tabular}

% % \subsection{Accents}
% \begin{tabular}{@{}l@{\ }l|l@{\ }l|l@{\ }l|l@{\ }l|l@{\ }l@{}}
% \`o   & \verb!\`o! &
% \'o   & \verb!\'o! &
% \^o   & \verb!\^o! &
% \~o   & \verb!\~o! &
% \=o   & \verb!\=o! \\
% \.o   & \verb!\.o! &
% \"o   & \verb!\"o! &
% \c o  & \verb!\c o! &
% \v o  & \verb!\v o! &
% \H o  & \verb!\H o! \\
% \c c  & \verb!\c c! &
% \d o  & \verb!\d o! &
% \b o  & \verb!\b o! &
% \t oo & \verb!\t oo! &
% \oe   & \verb!\oe! \\
% \OE   & \verb!\OE! &
% \ae   & \verb!\ae! &
% \AE   & \verb!\AE! &
% \aa   & \verb!\aa! &
% \AA   & \verb!\AA! \\
% \o    & \verb!\o! &
% \O    & \verb!\O! &
% \l    & \verb!\l! &
% \L    & \verb!\L! &
% \i    & \verb!\i! \\
% \j    & \verb!\j! &
% !`    & \verb!~`! &
% ?`    & \verb!?`! &
% \end{tabular}


% \subsection{Delimiters}
\begin{tabular}{@{}l@{\ }ll@{\ }ll@{\ }ll@{\ }ll@{\ }ll@{\ }l@{}}
`       & \verb!`!  &
``      & \verb!``! &
\{      & \verb!\{! &
\lbrack & \verb![! &
(       & \verb!(! &
\textless  &  \verb!\textless! \\
'       & \verb!'!  &
''      & \verb!''! &
\}      & \verb!\}! &
\rbrack & \verb!]! &
)       & \verb!)! &
\textgreater  &  \verb!\textgreater! \\
\end{tabular}



% \subsection{Line and page breaks}
% \settowidth{\MyLen}{\texttt{.pagebreak} }
% \begin{tabular}{@{}p{\the\MyLen}%
%                 @{}p{\linewidth-\the\MyLen}@{}}
% \verb!\\!          &  Begin new line without new paragraph.  \\
% \verb!\\*!         &  Prohibit pagebreak after linebreak. \\
% \verb!\kill!       &  Don't print current line. \\
% \verb!\pagebreak!  &  Start new page. \\
% \verb!\noindent!   &  Do not indent current line.
% \end{tabular}


% \subsection{Miscellaneous}
% \settowidth{\MyLen}{\texttt{.rule.w..h.} }
% \begin{tabular}{@{}p{\the\MyLen}%
%                 @{}p{\linewidth-\the\MyLen}@{}}
% \verb!\today!  &  \today. \\
% \verb!$\sim$!  &  Prints $\sim$ instead of \verb!\~{}!, which makes \~{}. \\
% \verb!~!       &  Space, disallow linebreak (\verb!W.J.~Clinton!). \\
% \verb!\@.!     &  Indicate that the \verb!.! ends a sentence when following
%                         an uppercase letter. \\
% \verb!\hspace{!$l$\verb!}! & Horizontal space of length $l$
%                                 (Ex: $l=\mathtt{20pt}$). \\
% \verb!\vspace{!$l$\verb!}! & Vertical space of length $l$. \\
% \verb!\rule{!$w$\verb!}{!$h$\verb!}! & Line of width $w$ and height $h$. \\
% \end{tabular}



\section{Tabular environments}

% \subsection{\texttt{tabbing} environment}
% \begin{tabular}{@{}l@{\hspace{1.5ex}}l@{\hspace{10ex}}l@{\hspace{1.5ex}}l@{}}
% \verb!\=!  &   Set tab stop. &
% \verb!\>!  &   Go to tab stop.
% \end{tabular}
% 
% Tab stops can be set on ``invisible'' lines with \verb!\kill!
% at the end of the line.  Normally \verb!\\! is used to separate lines.


\subsection{\texttt{tabular} environment}
% % \verb!\begin{array}[!\textit{pos}\verb!]{!\textit{Spalten}\verb!}!   \\
\verb!\begin{tabular}[!\textit{pos}\verb!]{!\textit{Spalten}\verb!}! \\
% \verb!\begin{tabular*}{!\textit{Breite}\verb!}[!\textit{pos}\verb!]{!\textit{Spalten}\verb!}!\
\verb!\begin{tabularx}{!\textit{Breite}\verb!}{!\textit{Spalten}\verb!}!\\
\verb!\begin{tabulary}{!\textit{Breite}\verb!}{!\textit{Spalten}\verb!}!      
  
tabularx: Erzeugt gleich-breite Spalten (X) mit Zeilenumbruch.\\
tabulary: Berechnet Spaltenbreite (R,L,C) automatisch und erlaubt Zeilenumbruch.


\subsubsection{Spaltentypen}
\settowidth{\MyLen}{\texttt{p}\{\textit{width}\} \ }
\begin{tabular}{@{}p{\the\MyLen}@{}p{\linewidth-\the\MyLen}@{}}
\texttt{l}    &   Linksbündig.  \\
\texttt{c}    &   Zentriert.  \\
\texttt{r}    &   Rechtsbündig. \\
\texttt{X}    &   tabularx: gleich-breite Spalten mit Zeilenumbruch. \\
\texttt{R}    &   tabulary: Rechtsbündig, Zeilenumbruch. \\ 
\texttt{L}    &   tabulary: Linksbündig, Zeilenumbruch. \\ 
\texttt{C}    &   tabulary: Zentriert, Zeilenumbruch. \\
% \verb!p{!\textit{width}\verb!}!  &  Spalte mit Zeilenumbruch, Breite muss
% angegeben werden!. 
% \\ \verb!@{!\textit{decl}\verb!}!   &  \noindent\Z{Insert
% \textit{decl} instead of inter-column space.} \\
% \verb!|!      &   \noindent\Z{Inserts a vertical line between
% columns.}
\end{tabular}

\begin{verbatim}
\begin{tabular}{lcr}
Das ist  & eine       & Tabelle
\\
mit      & 3 Spalten  & und 2 Zeilen
\\
\end{tabular}
\end{verbatim}

\begin{tabular}{lcr}
Das ist  & eine       & Tabelle
\\
mit      & 3 Spalten  & und 2 Zeilen
\\
\end{tabular}


% \subsubsection{\texttt{tabular} elements}
% \settowidth{\MyLen}{\texttt{.cline.x-y..}}
% \begin{tabular}{@{}p{\the\MyLen}@{}p{\linewidth-\the\MyLen}@{}}
% \verb!\hline!           &  Horizontal line between rows.  \\
% \verb!\cline{!$x$\verb!-!$y$\verb!}!  &
%                         Horizontal line across columns $x$ through $y$. \\
% \verb!\multicolumn{!\textit{n}\verb!}{!\textit{cols}\verb!}{!\textit{text}\verb!}! \\
%         &  A cell that spans \textit{n} columns, with \textit{cols} column specification.
% \end{tabular}


% \section{Math mode}
% To use math mode, surround text with \texttt{\$} or use
% \verb!\begin{equation}!.
% 
% \begin{tabular}{@{}l@{\hspace{1em}}l@{\hspace{2em}}l@{\hspace{1em}}l@{}}
% Superscript$^{x}$       &
% \verb!^{x}!             &  
% Subscript$_{x}$         &
% \verb!_{x}!             \\  
% $\frac{x}{y}$           &
% \verb!\frac{x}{y}!      &  
% $\sum_{k=1}^n$          &
% \verb!\sum_{k=1}^n!     \\  
% $\sqrt[n]{x}$           &
% \verb!\sqrt[n]{x}!      &  
% $\prod_{k=1}^n$         &
% \verb!\prod_{k=1}^n!    \\ 
% \end{tabular}
% 
% \subsection{Math-mode symbols}
% 
% % The ordering of these symbols is slightly odd.  This is because I had to put all the
% % long pieces of text in the same column (the right) for it all to fit properly.
% % Otherwise, it wouldn't be possible to fit four columns of symbols here.
% 
% \begin{tabular}{@{}l@{\hspace{1ex}}l@{\hspace{1em}}l@{\hspace{1ex}}l@{\hspace{1em}}l@{\hspace{1ex}} l@{\hspace{1em}}l@{\hspace{1ex}}l@{}}
% $\leq$          &  \verb!\leq!  &
% $\geq$          &  \verb!\geq!  &
% $\neq$          &  \verb!\neq!  &
% $\approx$       &  \verb!\approx!  \\
% $\times$        &  \verb!\times!  &
% $\div$          &  \verb!\div!  &
% $\pm$           & \verb!\pm!  &
% $\cdot$         &  \verb!\cdot!  \\
% $^{\circ}$      & \verb!^{\circ}! &
% $\circ$         &  \verb!\circ!  &
% $\prime$        & \verb!\prime!  &
% $\cdots$        &  \verb!\cdots!  \\
% $\infty$        & \verb!\infty!  &
% $\neg$          & \verb!\neg!  &
% $\wedge$        & \verb!\wedge!  &
% $\vee$          & \verb!\vee!  \\
% $\supset$       & \verb!\supset!  &
% $\forall$       & \verb!\forall!  &
% $\in$           & \verb!\in!  &
% $\rightarrow$   &  \verb!\rightarrow! \\
% $\subset$       & \verb!\subset!  &
% $\exists$       & \verb!\exists!  &
% $\notin$        & \verb!\notin!  &
% $\Rightarrow$   &  \verb!\Rightarrow! \\
% $\cup$          & \verb!\cup!  &
% $\cap$          & \verb!\cap!  &
% $\mid$          & \verb!\mid!  &
% $\Leftrightarrow$   &  \verb!\Leftrightarrow! \\
% $\dot a$        & \verb!\dot a!  &
% $\hat a$        & \verb!\hat a!  &
% $\bar a$        & \verb!\bar a!  &
% $\tilde a$      & \verb!\tilde a!  \\
% 
% $\alpha$        &  \verb!\alpha!  &
% $\beta$         &  \verb!\beta!  &
% $\gamma$        &  \verb!\gamma!  &
% $\delta$        &  \verb!\delta!  \\
% $\epsilon$      &  \verb!\epsilon!  &
% $\zeta$         &  \verb!\zeta!  &
% $\eta$          &  \verb!\eta!  &
% $\varepsilon$   &  \verb!\varepsilon!  \\
% $\theta$        &  \verb!\theta!  &
% $\iota$         &  \verb!\iota!  &
% $\kappa$        &  \verb!\kappa!  &
% $\vartheta$     &  \verb!\vartheta!  \\
% $\lambda$       &  \verb!\lambda!  &
% $\mu$           &  \verb!\mu!  &
% $\nu$           &  \verb!\nu!  &
% $\xi$           &  \verb!\xi!  \\
% $\pi$           &  \verb!\pi!  &
% $\rho$          &  \verb!\rho!  &
% $\sigma$        &  \verb!\sigma!  &
% $\tau$          &  \verb!\tau!  \\
% $\upsilon$      &  \verb!\upsilon!  &
% $\phi$          &  \verb!\phi!  &
% $\chi$          &  \verb!\chi!  &
% $\psi$          &  \verb!\psi!  \\
% $\omega$        &  \verb!\omega!  &
% $\Gamma$        &  \verb!\Gamma!  &
% $\Delta$        &  \verb!\Delta!  &
% $\Theta$        &  \verb!\Theta!  \\
% $\Lambda$       &  \verb!\Lambda!  &
% $\Xi$           &  \verb!\Xi!  &
% $\Pi$           &  \verb!\Pi!  &
% $\Sigma$        &  \verb!\Sigma!  \\
% $\Upsilon$      &  \verb!\Upsilon!  &
% $\Phi$          &  \verb!\Phi!  &
% $\Psi$          &  \verb!\Psi!  &
% $\Omega$        &  \verb!\Omega!  
% \end{tabular}
% \footnotesize
% 
% %\subsection{Special symbols}
% %\begin{tabular}{@{}ll@{}}
% %$^{\circ}$  &  \verb!^{\circ}! Ex: $22^{\circ}\mathrm{C}$: \verb!$22^{\circ}\mathrm{C}$!.
% %\end{tabular}
% 
% 
% 
% 
% \section{Sample \LaTeX\ document}
% \begin{lstlisting}
% \documentclass[11pt]{article}
% \usepackage{fullpage}
% \title{Template}
% \author{Name}
% \begin{document}
% \maketitle
% 
% \section{section}
% \subsection*{subsection without number}
% text \textbf{bold text} text. Some math: $2+2=5$
% \subsection{subsection}
% text \emph{emphasized text} text. \cite{WC:1953}
% discovered the structure of DNA.
% 
% A table:
% \begin{table}[!th]
% \begin{tabular}{|l|c|r|}
% \hline
% first  &  row  &  data \\
% second &  row  &  data \\
% \hline
% \end{tabular}
% \caption{This is the caption}
% \label{ex:table}
% \end{table}
% 
% The table is numbered \ref{ex:table}.
% \end{document}
% \end{lstlisting}



\rule{0.3\linewidth}{0.25pt}
\scriptsize

Copyright \copyright\ 2006 Winston Chang

% Should change this to be date of file, not current date.
\verb!$Revision: 1.13 $, $Date: 2008/05/29 06:11:56 $.!

http://www.stdout.org/$\sim$winston/latex/


\end{multicols}
\end{document}
